% ==================================================
% Caminho do template (pasta separada do projeto)
% Faz o LaTeX procurar a classe (.cls) e os pacotes (.sty)
% também dentro de modelo-cartaz/
% ==================================================
\makeatletter
\def\input@path{{./}{modelo-cartaz/}}
\makeatother

\documentclass[palette=basic,mode=light]{impaTeXPoster}

% Procura imagens primeiro na pasta do projeto, depois na do template
\graphicspath{{./}{images/}{modelo-cartaz/}}

\usepackage[utf8]{inputenc}
\usepackage[T1]{fontenc}
\usepackage{lipsum}

\title{Abrupt Changes in School Enrollment in Brazilian Basic Education:}
\subtitle{An Exploratory Analysis of INEP Data from 2011 to 2025}
\author{Daniela Moraes, Ingrid Carvalho}
\authoremail{al.daniela.moraes@impatech.edu.br, al.ingrid.carvalho@impatech.edu.br}
\posterlogo{images/impatech.png}    % logo à esquerda
%\rightposterlogo{images/tex.png}    % logo à direita

\begin{document}
\maketitle

\begin{columns}

% ==========================================
% LEFT COLUMN
% ==========================================
\column{0.5}

\begin{postercard}{Introduction}
Este é um exemplo de seção usando a classe ImpaTeXPoster.

\lipsum[1]
\end{postercard}

\begin{postercard}{Objetivo}
Mostrar como usar diferentes cores no ambiente.

\lipsum[2]
\end{postercard}

% ==========================================
% RIGHT COLUMN
% ==========================================
\column{0.5}

\begin{postercard}{Resultados}
\lipsum[3]
\end{postercard}

% primeiro parâmetro é a cor do Fundo do título
% segundo parâmetro é a cor do do testo
% terceiro parâmetro é cor do texto
% quarto parâmetro é a borda do postercard

\begin{postercard}{Metodologia}
\lipsum[4]
\end{postercard}

\end{columns}

\begin{columns}

%Left Column
\column{0.33}

\begin{postercard}{Metodologia}
\lipsum[5]
\end{postercard}

%Middle Column
\column{0.33}

\begin{postercard}{Metodologia}
\lipsum[6]
\end{postercard}


%Right Column
\column{0.33}

\begin{postercard}{Metodologia}
\lipsum[7]
\end{postercard}

\end{columns}


\posterfooter{https://github.com/impaTeX/modelo-cartaz/}{%

    \textbf{{\LARGE Referências:}}

    \begin{thebibliography}{99}
    \bibitem{popper}
    POPPER, Karl.
    \textit{A lógica da pesquisa científica}.
    São Paulo: Cultrix, 2002.

    \bibitem{chalmers}
    CHALMERS, A. F.
    \textit{O que é ciência, afinal?}
    São Paulo: Brasiliense, 1999.
    \end{thebibliography}
}

\end{document}